\cotefile{34}
\scanpage{18}
\mspage{16}
% archive sheet number: 17
% hand: Grothendieck
% sheet_kind: text
% element: text | top, line 1 | 4. Généralités
\deleted{\unsure{3}}4. \uline{Généralités sur les représentations $\ell$-adiques de $I$} $(\ell\ne p)$.

% element: text | upper, line 2 | Soit E un module
Soit $E$ un module de type fini sur $\mathbb Z_\ell$, donc
% element: display | upper, line 3 | E = lim E_nu
\[
E=\varprojlim_\nu \underbrace{E_\nu\deleted{\unsure{\otimes\mathbb Z/\ell^{\nu+1}\mathbb Z}}}_{E_\nu},\qquad E_\nu=E\otimes_{\mathbb Z_\ell}\mathbb Z/\ell^{\nu+1}\mathbb Z.
\tag*{(4.1)}
\]
% element: text | upper, line 4 | d'où
d'où
% element: display | upper, line 5 | G = Aut(E)
\[
G=\operatorname{Aut}_{\mathbb Z_\ell}(E)=\varprojlim_\nu\operatorname{Aut}_{\mathbb Z/\ell^\nu\mathbb Z}(E_\nu)
\tag*{(4.2)}
\]
% element: text | upper-middle, lines 6–8 | Par définition
Par définition, on munit $\operatorname{Aut}_{\mathbb Z_\ell}(E)$ de la structure de groupe topologique profini, limite projective des groupes finis apparaissant dans (4.1). Notons le

% element: text | middle, line 9 | Lemme 4.3
\uline{Lemme} 4.3 Le \deleted{\unsure{groupe}}\unsure{ss.-groupe} \deleted{\unsure{de}}
% element: display | middle, line 10 | G^1 = Ker
\[
G^1=\operatorname{Ker}\bigl(\operatorname{Aut}_{\mathbb Z_\ell}(E)\longrightarrow\operatorname{Aut}_{\mathbb Z/\ell\mathbb Z}(E_0)\bigr)
\]
% element: text | middle, lines 11–12 | est un sous-groupe
est un \unsure{sous-groupe|ss-groupe} ouvert de $G$, \deleted{\unsure{qui}} et c'est un pro-$\ell$-groupe.

% element: text | lower-middle, lines 13–15 | Dém.
\uline{Dém.} Le fait qu'il soit ouvert est trivial par définition. Que ce soit un pro-$\ell$-groupe \deleted{\unsure{est}} provient du fait que pour $\nu\geq 1$,
% element: display | lower-middle, line 16 | Ker(Aut(E_nu) to Aut(E_0))
\[
\operatorname{Ker}\bigl(\operatorname{Aut}(E_\nu)\longrightarrow\operatorname{Aut}(E_0)\bigr)
\]
% element: text | lower, lines 17–24 | est un l-groupe
est un $\ell$-groupe. Or si $u$ est tel que $u\equiv\operatorname{id}\;(\ell)$ \deleted{\unsure{i.e.} $(u-\unsure{1})(E_\nu)\subset\ell E_\nu$, on conclut par récurrence \unsure{que} $(u-1)^{\unsure{i}}\subset\ell^iE_\nu$, d'où pour $i=\nu+1$, $(u-1)^{\unsure{\nu}}E_\nu=0$, ce qui} alors $u$ induit l'identité sur le gradué associé à $E_\nu$ par la filtration (finie) \inserted{$\ell$-adique} (par les $\ell^iE_\nu$), d'où on conclut facilement par récurrence \inserted{sur $r$} que $u^{\ell^r}$ induit l'identité sur le
